\documentclass{article}
\usepackage{neurips_2026}

\usepackage[utf8]{inputenc}
\usepackage[T1]{fontenc}
\usepackage{hyperref}
\usepackage{url}
\usepackage{booktabs}
\usepackage{amsfonts}
\usepackage{microtype}
\usepackage{xcolor}
\usepackage{amsmath,amssymb,amsthm}
\usepackage{graphicx}
\usepackage{enumitem}

% ---------------------------------------------------------------- board machinery
% This file is a BOARD, not a draft. Status tags make the state visible at a glance.
\newcommand{\PROVED}{{\color{green!50!black}\textsc{[proved]}}}
\newcommand{\MEAS}{{\color{blue!70!black}\textsc{[measured]}}}
\newcommand{\OPEN}{{\color{orange!85!black}\textsc{[open]}}}
\newcommand{\DEAD}{{\color{red!75!black}\textsc{[retracted]}}}
\newcommand{\TODO}[1]{{\color{red!75!black}\textbf{TODO:} #1}}
\newcommand{\NOTE}[1]{{\color{gray}\small #1}}

\theoremstyle{plain}
\newtheorem{theorem}{Theorem}
\newtheorem{lemma}[theorem]{Lemma}
\newtheorem{proposition}[theorem]{Proposition}
\newtheorem{corollary}[theorem]{Corollary}
\theoremstyle{definition}
\newtheorem{definition}[theorem]{Definition}

\newcommand{\R}{\mathbb{R}}
\newcommand{\dW}{\Delta W}
\newcommand{\Jop}{J}
\newcommand{\Dstar}{\Delta_\star}
\newcommand{\Lam}{\Lambda}

\title{Operator-level optimisation, v2 --- working board}
\author{}
\date{}

\begin{document}
\maketitle

\NOTE{Numbers from \texttt{runs/theory/}; planning log \texttt{theory/13-paper-plan.md};
references in \texttt{references.bib} all read 2026-09-17.}

% =================================================================================
\section{Thesis and contributions}\label{sec:contrib}

\begin{quote}
\emph{In a gated network the ideal operator step is \textbf{not reachable}: the gates make the
contexts rank-deficient, so a fixed fraction of the operator gradient lies outside the reachable
set at every step, independently of depth. Within what \emph{is} reachable, gradient descent is
exactly \textbf{one Landweber iteration}, and it captures well under a third of the achievable
progress. Truncated Krylov iteration interpolates continuously between gradient descent and the
best achievable step. This gives the first direct control over the implicit bias of a
piecewise-linear network, and we use it to separate the part of the bias that is
\textbf{structural} (unreachable, set by the gates) from the part that is an artefact of
\textbf{solver conditioning} (correctable).}
\end{quote}

\begin{enumerate}[label=\textbf{C\arabic*.},leftmargin=*]
  \item \textbf{The ideal operator step is structurally unreachable in a ReLU network.} \MEAS{} \\
        $33$--$90\%$ of $\|\eta G\|$ lies outside $\mathrm{range}(M)$, \emph{even for a single
        input with $8\times$ more parameters than constraints}. Mechanism: every context is
        rank-limited by the gates. Not a depth effect --- a linear network at depth 128 realises
        the same step to $10^{-4}$. \S\ref{sec:unreach}.
  \item \textbf{Gradient descent is one Landweber step; the achievable ideal is the converged
        solve.} \PROVED{} \MEAS{} \\
        $\dW_\ell = -\mathrm{lr}\sum_x A_\ell^\top G(x) B_\ell^\top$ is exactly $M^\top$ applied
        to the target, so the induced step is $MM^\top\Dstar$ against $\Pi\Dstar =
        M(M^\top M)^+M^\top\Dstar$. Verified: at $k=1$, $\cos(\dW_{\mathrm{op}},-\nabla)=+1.0000$.
        \S\ref{sec:landweber}.
  \item \textbf{The dial.} \MEAS{} \\
        Truncating LSQR at $k$ iterations interpolates: $k=1$ is gradient descent, $k\to\infty$ is
        the achievable ideal. One principled parameter (Krylov regularisation), and cheap $k$ is
        what makes the method affordable. \S\ref{sec:dial}.
  \item \textbf{The bias decomposes into two mechanisms with different remedies.} \MEAS{} \\
        Unreachability (structural, gate rank) dominates at $L{=}2$; transfer-operator
        ill-conditioning (correctable) dominates at $L{=}4$, where $\mathrm{cond}(M)$ on its range
        rises $7.8\to1.7\!\times\!10^3$. \S\ref{sec:decomp}.
  \item \textbf{Method: minimal realisation, its certificate, and a constructive escape.}
        \PROVED{} \\
        Problem \eqref{eq:P}, the shared-multiplier condition \eqref{eq:kkt}, the residual $\rho$;
        the degeneracy theorem and the exact hierarchical recursion. \S\ref{sec:method}.
  \item \textbf{Conservation from the optimality form alone.} \PROVED{} \MEAS{} \\
        Any step of the form \eqref{eq:kkt} with a shared $\Lam$ conserves the balancedness
        invariant. Thm~\ref{thm:conserve}.
  \item \textbf{Real-data study.} \OPEN{} \\
        MNIST-1D running; full MNIST not yet feasible. \S\ref{sec:real}.
\end{enumerate}

\paragraph{Why not deep linear.} \DEAD{} Under balancedness the induced step
$T_W(G)=\sum_j(\Jop\Jop^\top)^{\frac{L-j}{L}}G(\Jop^\top \Jop)^{\frac{j-1}{L}}$ is a function of
$\Jop$ and $L$ alone \cite{arora2018optimization}, so the whole trajectory integrates in operator
space with no network and no solver: the weights are bookkeeping. Every comparison we ran there
reproduced known results \cite{saxe2014exact,gunasekar2017implicit,arora2019implicit,li2021towards}.
The deep linear case is the \emph{method's validation setting}, not its application.

% =================================================================================
\section{Setup}\label{sec:setup}

\begin{itemize}[leftmargin=*,itemsep=1pt]
  \item \textbf{Model.} ReLU MLP without biases, $f(x)=W_L\,\mathrm{relu}(W_{L-1}\cdots
        \mathrm{relu}(W_1x))$, hence $f(x)=\Jop(x)\,x$ with
        $\Jop(x)=W_LD_{L-1}(x)\cdots D_1(x)W_1$, $D_\ell(x)=\mathrm{diag}(\mathbb 1[z_\ell>0])$.
  \item \textbf{Contexts.} $\Jop(x)=A_\ell(x)W_\ell B_\ell(x)$ with
        $B_{\ell+1}=D_\ell W_\ell B_\ell$, $A_\ell=A_{\ell+1}W_{\ell+1}D_\ell$, $B_1=I$, $A_L=I$.
  \item \textbf{Per-sample operator gradient.} $G(x)=\partial\ell/\partial \Jop(x)$; rank one for
        squared loss ($(f-y)x^\top$) and for cross-entropy ($(\mathrm{softmax}(f)-e_y)x^\top$).
  \item \textbf{The reachable set.} $M(\dW) := \big(\sum_\ell A_\ell(x)\dW_\ell B_\ell(x)\big)_x$.
        Target $\Dstar=(-\eta G(x))_x$. \textbf{Achievable ideal} $\Pi\Dstar$, the projection of
        $\Dstar$ on $\mathrm{range}(M)$.
\end{itemize}

% =================================================================================
\section{The ideal step is unreachable}\label{sec:unreach}

\MEAS{} \texttt{studies/nonlinear.py}, $d{=}16$, exact least squares.

\begin{center}\small
\begin{tabular}{rrrrrr}
\toprule
$L$ & $n$ & unknowns & constraints & floor & generic \\
\midrule
2 & 1 & 512 & 256 & \textbf{0.375} & 0.000 \\
4 & 1 & 1024 & 256 & \textbf{0.501} & 0.000 \\
8 & 1 & 2048 & 256 & \textbf{0.328} & 0.000 \\
8 & 8 & 2048 & 2048 & 0.676 & 0.000 \\
8 & 32 & 2048 & 8192 & 0.616 & 0.866 \\
\bottomrule
\end{tabular}
\end{center}

\NOTE{The $n{=}1$ rows are the point: a single input, $8\times$ more parameters than constraints,
and a third to a half of the target is still unreachable. \emph{Mechanism:} every context is
rank-limited by the gates ($\mathrm{rank}\lesssim d/2$); for $\ell{=}1$, $B_1=I$, so the range
contributes $\mathrm{range}(A_1)\otimes\R^{d_{\mathrm{in}}}$ and the target's left factor
$f(x)-y$ generically leaves it. \emph{Contrast:} a deep linear network at depth 128 with the
operator collapsed to stable rank $1.1$ realises the same step to $1.7\times10^{-4}$.
\textbf{It is the gates, not the depth.}}
\TODO{Figure \texttt{images/fig-floor.pdf}: floor vs $n$, one line per $L$, with the generic
$\sqrt{1-k/m}$ baseline.}

% =================================================================================
\section{Gradient descent is one Landweber step}\label{sec:landweber}

\begin{proposition}\label{prop:landweber}\PROVED{}\MEAS{}
$\partial\mathcal L/\partial W_\ell = \sum_x A_\ell(x)^\top G(x)B_\ell(x)^\top$, i.e.\ the
coordinate gradient is $M^\top$ applied to the operator gradient. Hence
\[
  \text{GD's induced operator step} = M M^\top \Dstar,
  \qquad
  \text{achievable ideal} = \Pi\Dstar = M(M^\top M)^+M^\top\Dstar ,
\]
so the gap between them is exactly the conditioning of $MM^\top$ on its range.
\end{proposition}
\NOTE{Verified on MNIST-1D: the first LSQR iterate satisfies
$\cos(\dW_{\mathrm{op}},-\nabla)=+1.0000$. App.~\ref{app:landweber}.}

\paragraph{Consequence for our earlier measurements.} \DEAD{}
\texttt{theory/12} reported $\sin\to1$ between $\Delta \Jop$ and $G$ across a 108-run grid. That
compares against the \emph{ideal}, not the achievable ideal, and therefore forces $\sin\approx1$
whenever most of $G$ is unreachable, independently of what GD does.
\TODO{recompute the whole grid against $\Pi G$.}

% =================================================================================
\section{The dial}\label{sec:dial}

\MEAS{} \texttt{studies/dial.py}. Truncated LSQR on $M\dW=\Dstar$; $k$ is the iteration budget.
MNIST-1D, width 64, depth 4, batch 64:

\begin{center}\small
\begin{tabular}{rrr}
\toprule
$k$ & $\cos(\dW,-\nabla)$ & operator progress $\|M\dW\|/\|\Dstar\|$ \\
\midrule
1   & \textbf{+1.0000} & 0.248 \\
2   & +0.788 & 0.334 \\
5   & +0.515 & 0.491 \\
15  & +0.332 & 0.646 \\
50  & +0.169 & 0.776 \\
150 & +0.087 & \textbf{0.839} \\
\bottomrule
\end{tabular}
\end{center}

\NOTE{Exactly calibrated at both ends. \textbf{Gradient descent realises $24.8\%$ of the ideal
operator step against an achievable ceiling of $83.9\%$} --- under a third of what the gates
permit.}
\TODO{Figure \texttt{images/fig-dial.pdf}: the two curves against $k$.}

% =================================================================================
\section{Decomposition of the bias}\label{sec:decomp}

\MEAS{} \texttt{studies/achievable.py}, ReLU MLP, $d{=}16$.

\begin{center}\small
\begin{tabular}{rrrrrr}
\toprule
$L$ & $n$ & achievable frac. & $\cos(\mathrm{GD},\text{ideal})$ & $\cos(\mathrm{GD},\text{achievable})$ & $\mathrm{cond}(M)$ \\
\midrule
2 & 1  & 0.864 & 0.736 & \textbf{0.852} & 5.4 \\
2 & 16 & 0.448 & 0.367 & \textbf{0.820} & 7.8 \\
4 & 1  & 0.631 & 0.311 & \textbf{0.493} & $1.3\!\times\!10^{2}$ \\
4 & 16 & 0.752 & 0.418 & \textbf{0.556} & $1.7\!\times\!10^{3}$ \\
\bottomrule
\end{tabular}
\end{center}

\begin{enumerate}[leftmargin=*,itemsep=1pt]
  \item \textbf{Unreachability} (gate rank). Not correctable. Dominant at $L{=}2$: correcting for
        it moves GD from $0.367$ to $0.820$.
  \item \textbf{Ill-conditioning of the transfer operator.} Correctable --- this is what the dial
        buys. Dominant at $L{=}4$ ($0.418\to0.556$ only) as $\mathrm{cond}$ climbs three orders.
\end{enumerate}
\NOTE{The balance shifts toward the correctable mechanism with depth. This is the paper's central
measurement.}
\TODO{Figure \texttt{images/fig-decomposition.pdf}. Extend to $L\in\{8,16\}$ and to CReLU.}

% =================================================================================
\section{Method}\label{sec:method}

\begin{definition}[(P)]
\begin{equation}\label{eq:P}
  \min_{\dW}\ \sum_\ell\|\dW_\ell\|_F^2 \quad\text{s.t.}\quad \prod_\ell(W_\ell+\dW_\ell)=P_\star .
\end{equation}
\end{definition}

\begin{lemma}[shared multiplier]\label{lem:kkt}\PROVED{}
At a regular optimum there is a single $\Lam$ with
\begin{equation}\label{eq:kkt}
  \dW_\ell = A'^\top_\ell\,\Lam\,B'^\top_\ell\quad\text{for every }\ell,
\end{equation}
contexts taken at the \emph{updated} weights. The certificate
$\rho(\dW)=\min_\Lam\sum_\ell\|\dW_\ell-A'^\top_\ell\Lam B'^\top_\ell\|^2/\sum_\ell\|\dW_\ell\|^2$
tests any solver's output without knowing the optimum.
\end{lemma}
\NOTE{GD satisfies \eqref{eq:kkt} for the \emph{linearised} constraint with contexts at the
\emph{current} weights. Equation \eqref{eq:kkt} is self-referential, which is why a fixed-point
iteration from $\dW=0$ never leaves and the escape must be constructive.
\textbf{Caveat:} necessary only at a regular optimum, and the constraint Jacobian loses rank
exactly where the operator collapses. \TODO{characterise LICQ for (P).}}

\begin{theorem}[degeneracy]\label{thm:degen}\PROVED{}\MEAS{}
At $W=0$ the product map is homogeneous of degree $L$, so all derivatives of order $<L$ vanish and
every null-consistent derivative-based rule of order $p<L$ returns $\dW=0$ --- including natural
gradient \cite{bernacchia2018exact} and Layerwise LQR \cite{dufortlabbe2026layerwise}, and
including a first-order solver applied to the exact objective.
\end{theorem}

\begin{theorem}[exact recursion]\label{thm:exact}\PROVED{}\MEAS{}
From a degenerate base point the hierarchical recursion solves \eqref{eq:P} exactly for every $L$
iff each split uses the depth-proportional exponent $t=a/(a+b)$; the symmetric split is exact iff
$L$ is a power of two. Verified to $1.000000$ for $L=2\ldots128$ at $d{=}32$, in closed form,
${<}0.05$\,s.
\end{theorem}

\begin{theorem}[conservation]\label{thm:conserve}\PROVED{}\MEAS{}
Any flow $\dot W_\ell=A_\ell^\top\Lam B_\ell^\top$ with a \emph{shared} $\Lam$ conserves
$W_{\ell+1}^\top W_{\ell+1}-W_\ell W_\ell^\top$. So the Du--Hu--Lee / Arora balancedness invariant
\cite{du2018algorithmic,arora2018optimization} is a property of minimal realisations, not of
gradient descent.
\end{theorem}
\NOTE{Measured: discrepancy $1.5\times10^{-15}$ for random shared $\Lam$ vs $1.4\times10^{1}$ for
an independent per-layer step.}

% =================================================================================
\section{Real data}\label{sec:real}

\OPEN{} \texttt{studies/dial.py}. MNIST-1D ($40\!\to\!10$, cross-entropy), ReLU MLP width 64,
depth 4, 1000 train / 500 test, batch 64, 400 steps. Every arm gets a hyper-parameter grid.

\NOTE{Partial, seed 0: the $k$-dial is monotone in both train and \emph{test} loss ---
$k{=}1$: train 1.973 / test 2.077; $k{=}5$: 1.418 / 1.906; $k{=}15$: 1.068 / 1.826;
$k{=}50$: 0.847 / 1.816. Baselines reach far lower train loss and far \emph{higher} test loss
(GD $\mathrm{lr}{=}0.2$: 0.241 / 2.688; Adam $\mathrm{lr}{=}0.01$: 0.063 / 4.454).}

\TODO{\textbf{Fairness defect, mine:} the baselines get an lr sweep while the operator arm has
$\eta$ fixed at $1.0$. That is the same asymmetry that sank the previous submission, in reverse.
Sweep $\eta$, and compare \emph{at matched train loss} --- the arms currently sit at very
different points on the fitting curve, so the test-loss gap is confounded with early stopping.}
\TODO{full MNIST ($784\!\to\!10$) is not feasible as written: a matvec is ${\sim}500$\,MFLOP at
batch 128, so $300$ LSQR iterations is ${\sim}150$\,s \emph{per step}. Options: harder
minibatching, small $k$, downsampling, or restructuring the solve.}

% =================================================================================
\section{Positioning, and what is not claimed}\label{sec:notclaimed}

\begin{itemize}[leftmargin=*,itemsep=1pt]
  \item \DEAD{} No new characterisation of the deep linear implicit bias
        (\cite{arora2018optimization} Thm~1); no missing closed form for fixed-gate networks
        (absorbing gates into $M_\ell=D_\ell W_\ell D_{\ell-1}$ makes the FGLN flow literally the
        linear flow, so \cite{haas2026why} Prop.~4.5 is correct); $K_\ell\ge L$ is not
        hypothesis-free (only $K_\ell\ge L(s_\ell/g)^{2/L}$); no universal benefit from removing
        the bias --- on matrix sensing the \emph{biased} arm recovers the low-rank truth.
  \item \textbf{Not an optimiser paper.} The method realises at best ${\sim}80\%$ of the step it
        targets, at $10\times$ the wall clock. Benchmarking is the wrong frame; it is an
        instrument. \emph{Training to comparable loss is a validity precondition, not a claim.}
  \item \textbf{Prior art to position against.} $2-2/L$ by AM--GM \cite{rawal2026saddle};
        Riemannian metric view \cite{bah2022learning}; depth removal by constrained
        parameterisation \cite{brechet2023deep}; minimal perturbation for an output change
        \cite{evans2026theory}; layer-wise targets \cite{frerix2018proximal}; gated linear
        implicit bias, asymptotic \cite{gatedlinear2022implicit}.
\end{itemize}

% =================================================================================
\section{Open items, ordered}\label{sec:open}

\begin{enumerate}[leftmargin=*,itemsep=1pt]
  \item Fix the $\eta$ sweep and the matched-train-loss comparison in \S\ref{sec:real}.
  \item Recompute the 108-run grid against $\Pi G$ instead of $G$.
  \item Extend \S\ref{sec:decomp} to $L\in\{8,16\}$, to CReLU, and to MNIST-1D.
  \item Make full MNIST feasible, or state why it is not and drop it.
  \item Cost-to-go at a general base point (closes Thm~\ref{thm:exact} off the degenerate point);
        currently only the aligned/separated-spectrum case, ratio $1.18$--$1.36$.
  \item LICQ for (P).
\end{enumerate}

% =================================================================================
\bibliographystyle{plain}
\bibliography{references}

\appendix
\section{Proofs}
\subsection{Lemma \ref{lem:kkt}}\label{app:kkt}\TODO{three lines.}
\subsection{Proposition \ref{prop:landweber}}\label{app:landweber}\TODO{$\partial_{W_\ell}$ of the
per-sample loss is $A_\ell^\top G B_\ell^\top$; stack over $\ell$ and $x$.}
\subsection{Theorem \ref{thm:degen}}\label{app:degen}\TODO{homogeneity; define null-consistency.}
\subsection{Theorem \ref{thm:exact}}\label{app:exact}\TODO{Bellman with
$V_a(Y)=a\sum_i\sigma_i(Y)^{2/a}$; AM--GM; power-of-two characterisation.}
\subsection{Theorem \ref{thm:conserve}}\label{app:conserve}\TODO{$\dot W_{\ell+1}^\top W_{\ell+1}$
and $\dot W_\ell W_\ell^\top$ are transposes.}

\section{Additional measurements}
\NOTE{ALS is not a minimal realisation ($\|\dW\|$ up to $2413\times$ the minimum, $\rho\approx1$)
--- appendix only, it explains the warm-start dependence of the previous submission. Unbalanced
excess $K_\ell/L$: exactly $1.000$ when balanced, $38$--$230$ at Xavier depth 128. Reachability
floor $\le1.7\times10^{-4}$ for deep linear at depth 128. The 108-run grid,
\texttt{theory/12-grid.md}, pending recomputation against $\Pi G$.}

\end{document}
